%% evolutionary-loop.tex: the agentic evolutionary loop (Chapter 10).
%% House conceptual style; Mutate is drawn as a node distinct from Generate, since it
%% is the one operator absent from the generate-evaluate-select loop of Chapter 9.
%% _concept-style.tex: shared TikZ styles for the course's CONCEPTUAL block diagrams.
%% The leading underscore marks this as the one file in figures/ that is not a
%% figure; every other figures/*.tex holds exactly one figure body.
%%
%% \input this at the top of every conceptual figure (before \begin{tikzpicture})
%% so each such diagram speaks one visual language: rounded "block" nodes, stealth
%% arrows, and natural `<hue>!15' fills. Redefining these styles on each \input is
%% idempotent and harmless. The reference for the house parameters is
%% resources/STYLE-figures.md. Figure~\ref{figure:AgentLoop} is the exemplar.
%%
%% Scope: use for conceptual/relational diagrams (boxes-and-arrows). Do NOT use for
%% product mock-ups (the rig figures have their own shared language) or for
%% data/plot figures.
\tikzset{
  % the standard block node: geometry only — apply a `fill=<hue>!15' per node.
  conceptbox/.style={draw, rounded corners=4pt, minimum width=2.4cm,
                     minimum height=0.85cm, align=center, font=\small},
  % a flow arrow between blocks.
  conceptflow/.style={->, >=stealth', thick},
  % an edge/annotation label.
  conceptlabel/.style={font=\scriptsize, align=center},
}%

\begin{tikzpicture}
% One-time seed on the left; the steady-state cycle runs clockwise on the right.
\node[conceptbox, fill=yellow!20, minimum width=3.4cm] (pop) at (0, 2.0) {Initial population $\Population_0$};
\node[conceptbox, fill=blue!8]    (gen) at (0,    0)    {Generate};
\node[conceptbox, fill=green!20]  (fit) at (5.2,  0)    {Evaluate\\(Fitness)};
\node[conceptbox, fill=orange!20] (sel) at (10.4, 0)    {Select};
\node[conceptbox, fill=blue!15]   (mut) at (5.2, -2.2)  {Mutate};

\draw[conceptflow] (pop) -- node[conceptlabel, midway, right]{seed}    (gen);
\draw[conceptflow] (gen) -- node[conceptlabel, midway, above]{initial} (fit);
\draw[conceptflow] (fit) -- node[conceptlabel, midway, above]{scored}  (sel);
\draw[conceptflow] (sel.south) -- (sel.south |- mut.east)
    -- node[conceptlabel, midway, below]{survivors $\Population_{t+1}$} (mut.east);
\draw[conceptflow] (mut) -- node[conceptlabel, midway, right]{candidates} (fit);
\end{tikzpicture}%
